% Trustworthy RAG final-report scaffold.
% Compile after replacing every [TBD] item and adding the locked raw-B2 result.
\documentclass[conference]{IEEEtran}
\usepackage[T1]{fontenc}
\usepackage[utf8]{inputenc}
\usepackage{booktabs}
\usepackage{amsmath}
\usepackage{graphicx}
\usepackage{xcolor}
\usepackage{url}
\newcommand{\todo}[1]{\textcolor{red}{[TBD: #1]}}

\title{Risk-Aware Retrieval for Defending Retrieval-Augmented Generation Against Indirect Prompt Injection}
\author{\IEEEauthorblockN{\todo{Student 1 name} \and \todo{Student 2 name}}\IEEEauthorblockA{\todo{Course, University of Auckland}\\\todo{student email addresses}}}

\begin{document}
\maketitle

\begin{abstract}
Retrieval-augmented generation (RAG) improves knowledge access, but it also
creates a trust-boundary problem: a retrieved document may contain instructions
that compete with the user's request. This project studies whether passage
selection can account for both relevance and prompt-injection risk before
retrieved text reaches a generator. We implement a reproducible EmailQA
pipeline around BIPIA, dense bge-m3 retrieval, the PIGuard injection detector,
and a fixed Qwen generator. We compare unprotected RAG, PIGuard hard filtering,
a stricter risk-only rule, and a relevance--risk selection rule. The primary
protocol uses a frozen, stratified 300-case attacked test set plus 20 clean
cases; it reports attack success, malicious-context inclusion, answer-bearing
evidence retention, clean-task behaviour, abstention, and confidence intervals.
Preliminary runs show that filtering sharply reduces attack-context inclusion
while potentially discarding useful evidence. The final table will insert the
locked original-PIGuard baseline and audit results. Rather than claiming
universal robustness, the study identifies the security--utility trade-off
created by detector-driven retrieval decisions.
\end{abstract}

\begin{IEEEkeywords}
retrieval-augmented generation, indirect prompt injection, trustworthy AI,
prompt-injection detection, retrieval security
\end{IEEEkeywords}

\section{Introduction}
RAG systems routinely place retrieved natural-language documents beside a
user's request. This improves factual grounding, but it also mixes data and
instructions in one model input. Indirect prompt injection exploits that
boundary: an attacker embeds an instruction in external content, and the model
may follow it after retrieval \cite{greshake2023}. The resulting failure can
redirect a task even if the user never supplied the malicious instruction.

Our research question is: \emph{Can a RAG system estimate the trustworthiness
of retrieved passages before generation and reduce indirect prompt injection
without unnecessarily losing answer-bearing evidence?} We test the hypothesis
that using relevance alongside an injection-risk score can preserve more useful
content than a detector-only hard filter at similar attack-context removal.

This work makes three course-project contributions. First, it implements a
versioned BIPIA EmailQA experimental pipeline with deterministic manifests,
caching, and an explicit retrieval pool. Second, it compares an original
PIGuard chunk-level filter with a windowed detector ablation, a risk-only rule,
and a relevance--risk selector. Third, it evaluates security and utility
together, and records limitations rather than treating a low attack-success
rate as sufficient evidence of general robustness.

\section{Related Work}
RAG couples a retriever and generator \cite{lewis2020,karpukhin2020}, but the
retrieved corpus can be adversarial. Greshake \emph{et al.} demonstrated
indirect prompt injection in LLM-integrated applications \cite{greshake2023}.
Subsequent benchmarks formalise attacks and defences \cite{liu2024,bipia2025},
and security benchmarks such as SafeRAG examine broader RAG failure modes
\cite{saferag2025}. PIGuard focuses on reducing detector over-defence caused
by superficial trigger words \cite{piguard2025}. Our project does not train a
new generator or detector. Instead, it asks how a detector score should alter
retrieval selection, with explicit evidence-retention measurements.

\section{Methodology}
\subsection{Threat model and retrieval setting}
The attacker controls an instruction inserted into an EmailQA document, but
does not control the user's question, model weights, or scoring code. Each
query retrieves from a small constructed pool: the attacked target email and
four seeded distractors that do not mention the query entity. Hence our claims
are limited to this BIPIA EmailQA RAG setting, not to arbitrary web-scale RAG.

\subsection{Methods}
Let $q$ be a question, $c$ a sentence-level candidate chunk,
$\operatorname{rel}(q,c)$ dense cosine relevance, $r(c)$ PIGuard injection
risk, and $\operatorname{conflict}(q,c)$ an optional context-conflict value.
Our selection score is
\begin{equation}
s(q,c)=\alpha\operatorname{rel}(q,c)-\beta r(c)-\gamma\operatorname{conflict}(q,c).
\end{equation}
Chunks meeting the frozen keep rule are returned in source-document order.

\begin{table}[t]
\caption{Compared methods. Raw B2 is the existing PIGuard baseline; B2-win
is a project window-scoring ablation.}
\label{tab:methods}
\centering
\begin{tabular}{ll}
\toprule
Method & Selection rule \\
\midrule
B0 & Relevance-only RAG; no defence \\
B2 (raw) & Chunk-level PIGuard score; remove risk $\geq0.5$ \\
B2-win & Windowed PIGuard score; remove risk $\geq0.5$ \\
B3 & Windowed risk-only threshold ($r\leq0.25$) \\
Ours & Relevance--risk--conflict score with frozen keep rule \\
\bottomrule
\end{tabular}
\end{table}

\subsection{Protocol and measures}
The primary split is a deterministic stratified sample of 300 BIPIA test
attacks across 15 attack families and three insertion positions, plus 20 clean
emails (10 known-answer and 10 unknown-answer). The generator is
Qwen2.5-VL-32B-Instruct at temperature zero; bge-m3 supplies embeddings;
PIGuard scores risk locally. The official BIPIA attack evaluator is used,
with Qwen2.5-VL-72B-AWQ answering model-judged criteria in place of the
original GPT-3.5. UNKNOWN judgement outputs count as not successful and are
reported.

Primary security measures are official attack success rate (ASR), attack text
in the final context, and attack removal. Utility measures are answer-bearing
chunk retention on the 120 answerable attacked cases, clean-case accuracy on
20 cases, clean chunks removed, and abstention over all 320 cases. We report
rates with Wilson 95\% intervals and paired exact comparisons. A 100-case
pilot is infrastructure-only and not pooled with the Main Test; a disjoint
45-case follow-up is supplementary and not used for tuning.

\section{Implementation and Results}
The implementation records model revisions, manifests, selection decisions,
generated outputs, and judge outputs, while caching deterministic intermediate
steps. It uses sentence-level defence units within whole-email documents. The
experiment is reproducible from the public repository\footnote{\todo{Insert
public GitHub URL after final review.}} without committing credentials or model
weights.

Table~\ref{tab:results} is intentionally incomplete until the raw B2 Main-Test
run and human-audit sensitivity check are locked. The already-completed
windowed comparisons should not be relabelled as original B2. Current evidence
supports a cautious interpretation: all evaluated defences reduced ASR relative
to B0, while Ours retained more answer-bearing chunks than the two stricter
windowed rules. It does \emph{not} support a claim that the conflict term adds
value, because it did not change frozen selections.

\begin{table*}[t]
\caption{Primary Main-Test result template. Fill only from the locked metric
report; do not pool pilot or follow-up cases.}
\label{tab:results}
\centering
\begin{tabular}{lccccc}
\toprule
Method & ASR ($n=300$) & Attack in context & Removal & Evidence retention ($n=120$) & Clean ($n=20$) \\
\midrule
B0 & \todo{locked} & \todo{locked} & \todo{locked} & \todo{locked} & \todo{locked} \\
B2 (raw) & \todo{pending raw-B2 run} & \todo{pending} & \todo{pending} & \todo{pending} & \todo{pending} \\
B2-win & \todo{locked} & \todo{locked} & \todo{locked} & \todo{locked} & \todo{locked} \\
B3 & \todo{locked} & \todo{locked} & \todo{locked} & \todo{locked} & \todo{locked} \\
Ours & \todo{locked} & \todo{locked} & \todo{locked} & \todo{locked} & \todo{locked} \\
\bottomrule
\end{tabular}
\end{table*}

\subsection{Limitations}
This evaluation uses one generator and a small clean set, so it cannot estimate
small utility differences precisely. The evaluator model differs from BIPIA's
original setup; model-judged UNKNOWN cases are a source of uncertainty. The
constructed retrieval pool and fixed temperature also limit generalisation.
Finally, the 45-case post-main template comparison is evidence of robustness
checking, not an additional primary test set.

\section{Conclusion and Future Work}
This project operationalises trustworthiness as a retrieval-time selection
problem rather than only a generator prompt problem. The final locked
comparison will determine whether relevance-aware selection improves evidence
retention at comparable security in the studied BIPIA setting. Future work
should evaluate more generators, larger clean sets, adaptive attackers,
web-scale corpora, and a conflict signal that is independently validated before
being included in a scoring rule.

\section*{Contribution Breakdown}
\todo{Replace with a factual breakdown of each group member's work, including
literature review, implementation, experiment verification, writing, slides,
and recording.}

\begin{thebibliography}{00}
\bibitem{lewis2020} P. Lewis \emph{et al.}, ``Retrieval-augmented generation for knowledge-intensive NLP tasks,'' in \emph{Advances in Neural Information Processing Systems}, vol. 33, 2020.
\bibitem{karpukhin2020} V. Karpukhin \emph{et al.}, ``Dense passage retrieval for open-domain question answering,'' in \emph{Proc. EMNLP}, 2020.
\bibitem{greshake2023} K. Greshake \emph{et al.}, ``Not what you've signed up for: Compromising real-world LLM-integrated applications with indirect prompt injection,'' in \emph{Proc. AISec}, 2023.
\bibitem{liu2024} Y. Liu \emph{et al.}, ``Formalizing and benchmarking prompt injection attacks and defenses,'' in \emph{Proc. USENIX Security}, 2024.
\bibitem{bipia2025} J. Yi \emph{et al.}, ``Benchmarking and defending against indirect prompt injection attacks on large language models,'' in \emph{Proc. KDD}, 2025.
\bibitem{saferag2025} X. Liang \emph{et al.}, ``SafeRAG: Benchmarking security in retrieval-augmented generation of large language model,'' in \emph{Proc. ACL}, 2025.
\bibitem{piguard2025} H. Li, X. Liu, N. Zhang, and C. Xiao, ``PIGuard: Prompt injection guardrail via mitigating overdefense for free,'' in \emph{Proc. ACL}, 2025.
\bibitem{agentdojo2024} E. Debenedetti \emph{et al.}, ``AgentDojo: A dynamic environment to evaluate prompt injection attacks and defenses for LLM agents,'' in \emph{Advances in Neural Information Processing Systems}, 2024.
\bibitem{poisonedrag2025} W. Zou \emph{et al.}, ``PoisonedRAG: Knowledge corruption attacks to retrieval-augmented generation of large language models,'' in \emph{Proc. USENIX Security}, 2025.
\bibitem{zhan2025} Q. Zhan \emph{et al.}, ``Adaptive attacks break defenses against indirect prompt injection attacks on LLM agents,'' in \emph{Findings of NAACL}, 2025.
\end{thebibliography}
\end{document}
